% Tier 1 parity doc: CJK (Chinese) — native-word nodes + locale linebreak.
% Exercises do_locale_linebreaks and inter-character glue for CJK.
\documentclass[12pt]{article}
\usepackage{fontspec}
\setmainfont{LibertinusSerif-Regular.otf}
\newfontfamily\cjkfont{FandolSong-Regular.otf}[Script=CJK]
\begin{document}
\section*{CJK line breaking}

English lead-in, then a block of Chinese text set with a CJK font to
drive the locale-aware line breaker over a long run of wide glyphs.

{\cjkfont
排版系统需要正确处理中文文本的断行与字间距。这是一段足够长的中文，
用来测试在没有空格的情况下，断行算法能否在字符之间正确地折行，并且
在整段范围内保持一致的字间距与标点处理。数字一二三四五六七八九十。
}

Back to English to confirm the surrounding paragraph still typesets.

\end{document}
